\subsection{卷积与线性表示}

命题 11. --- 设 G 为局部紧群，E 为拟完备局部凸空间，U 为 G 在 E 中的连续表示。

\begin{enumerate}
\def\labelenumi{(\roman{enumi})}
\item
  若 \(\lambda \in \mathcal{C}'(\mathbf{G})\)，\(\mu \in \mathcal{C}'(\mathbf{G})\)，则 \(U(\lambda * \mu) = U(\lambda)U(\mu)\)。
\item
  假设 \(\mathbf{E}\) 是 Banach 空间，并置 \(\rho(s) = \|U(s)\|\)（\(s\in\mathbf{G}\)）。若 \(\lambda\in\mathcal{M}^{\rho}(\mathbf{G})\)，\(\mu\in\mathcal{M}^{\rho}(\mathbf{G})\)，则 \(U(\lambda*\mu)=U(\lambda)U(\mu)\)。
\end{enumerate}

设 \(\lambda, \mu\) 属于 \(\mathcal{C}'(\mathrm{G})\)。对任意 \(x \in \mathrm{E}\)，应用第六章 §1, No.~1 的命题 1 和命题 4，即得

\[
\begin{array}{r l} & U (\lambda * \mu) x = \int_ {\mathbf {G}} U (s) x d (\lambda * \mu) (s) \\ & \quad = \int_ {\mathbf {G} \times \mathbf {G}} U (s t) x d \lambda (s) d \mu (t) = \int_ {\mathbf {G} \times \mathbf {G}} U (s) U (t) x d \lambda (s) d \mu (t) \\ & \quad = \int_ {\mathbf {G}} U (\lambda) U (t) x d \mu (t) = U (\lambda) \int_ {\mathbf {G}} U (t) x d \mu (t) = U (\lambda) U (\mu) x, \end{array}
\]

由此得 (i)。类似的论证可应用于情形 (ii)。

仍设 G 为局部紧群，并假设 G 连续地左作用于局部紧空间 X。这就定义了（§2, No.~4）G 的一个连续线性表示 \(\pmb{\gamma}\)，它作用在 \(\mathcal{M}(\mathrm{X})\) 中（后者赋予在 \(\mathcal{K}(\mathrm{X})\) 中的紧收敛拓扑）。

命题 12. --- 若 \(\lambda \in \mathcal{C}'(\mathbf{G})\) 且 \(\mu \in \mathcal{M}(\mathbf{X})\)，则

\[
\gamma (\lambda) \mu = \lambda * \mu .
\]

由 §1, No.~5 的命题 7，

\[
\lambda * \mu = \int_ {G} (\varepsilon_ {s} * \mu) d \lambda (s).
\]

而 \(\varepsilon_{s}*\mu = \pmb {\gamma}(s)\mu\)（No.~2 公式 (5)），且

\[
\int_ {\mathsf {G}} \left(\pmb {\gamma} (s) \mu\right) d \lambda (s) = \pmb {\gamma} (\lambda) \mu
\]

这是按 \(\pmb{\gamma}(\lambda)\) 的定义得到的。

推论. --- 映射 \((\lambda,\mu)\mapsto\lambda*\mu\)（从 \(\mathcal{C}^{\prime}(\mathrm{G})\times\mathcal{M}(\mathrm{X})\) 到 \(\mathcal{M}(\mathrm{X})\) 中）相对于 \(\mathcal{C}^{\prime}(\mathrm{G})\) 的等度连续子集与 \(\mathcal{M}(\mathrm{X})\) 的紧子集是亚连续的（\(\mathcal{C}^{\prime}(\mathrm{G})\) 与 \(\mathcal{M}(\mathrm{X})\) 分别赋予在 \(\mathcal{C}(\mathrm{G})\) 与 \(\mathcal{H}(\mathrm{X})\) 中的紧收敛拓扑）。

因为 \(\mathcal{M}(\mathrm{X})\) 赋予在 \(\mathcal{K}(\mathrm{X})\) 中的紧收敛拓扑后是拟完备的，所以映射 \((\lambda, \mu) \mapsto \pmb{\gamma}(\lambda)\mu\)（从 \(\mathcal{C}'(\mathrm{G}) \times \mathcal{M}(\mathrm{X})\) 到 \(\mathcal{M}(\mathrm{X})\) 中）相对于 \(\mathcal{C}'(\mathrm{G})\) 的等度连续子集与 \(\mathcal{M}(\mathrm{X})\) 的紧子集是亚连续的（§2, No.~6）。于是只需应用命题 12。

备注. --- 1) 设 \(\lambda_0 \in \mathcal{C}'(\mathbf{G})\)。映射 \(\mu \mapsto \lambda_0 * \mu\)（从 \(\mathcal{M}(\mathbf{X})\) 到 \(\mathcal{M}(\mathbf{X})\) 中）是淡连续的。因为，设 \(f \in \mathcal{K}(\mathbf{X})\)。我们有

\[
\langle \lambda_ {0} * \mu , f \rangle = \int f (s x) d \lambda_ {0} (s) d \mu (x) = \langle \mu , g \rangle ,
\]

其中 \(g(x) = \int f(sx) d\lambda_0(s)\)。而 \(g\) 是连续的（第七章 §1, No.~1 引理 1）。另一方面，设 \(S\) 为 \(\lambda_0\) 的支集，\(K\) 为 \(f\) 的支集。条件 \(sx \in K\) 与 \(s \in S\) 蕴含 \(x \in S^{-1}K\)；因此 \(g\) 的支集包含于 \(S^{-1}K\)，从而 \(g \in \mathcal{K}(X)\)。于是 \(\langle \lambda_0 * \mu, f \rangle = \langle \mu, g \rangle\) 是 \(\mu\) 的淡连续函数，这就证明了我们的断言。

\begin{enumerate}
\def\labelenumi{\arabic{enumi})}
\setcounter{enumi}{1}
\tightlist
\item
  设 \(\mu_0 \in \mathcal{M}(\mathrm{X})\)。映射 \(\lambda \mapsto \lambda * \mu_0\)（从 \(\mathcal{C}'(\mathrm{G})\) 到 \(\mathcal{M}(\mathrm{X})\) 中）对于拓扑 \(\sigma(\mathcal{C}'(\mathrm{G}), \mathcal{C}(\mathrm{G}))\) 与 \(\sigma(\mathcal{M}(\mathrm{X}), \mathcal{K}(\mathrm{X}))\) 是连续的。因为，设 \(f \in \mathcal{K}(\mathrm{X})\)。置 \(h(s) = \int f(sx) d\mu_0(x)\)，则 \(\langle f, \lambda * \mu_0 \rangle = \langle h, \lambda \rangle\)，且 \(h \in \mathcal{C}(\mathrm{G})\)（第七章 §1, No.~1 引理 1）。
\end{enumerate}

命题 13. --- 映射 \((s,\mu)\mapsto\boldsymbol{\gamma}(s)\mu\)（从 \(\mathbf{G}\times\mathcal{M}_{+}(\mathbf{X})\) 到 \(\mathcal{M}_{+}(\mathbf{X})\) 中）当 X 上的正测度所成的集合 \(\mathcal{M}_{+}(\mathbf{X})\) 赋予淡拓扑时是连续的。

由于 \(\pmb{\gamma}(s)\mu = \pmb{\gamma}(ss_0^{-1})\pmb{\gamma}(s_0)\mu\)，根据备注 1，只需证明所考虑的映射在形如 \((e,\mu_0)\)（其中 \(\mu_0\in \mathcal{M}_+(\mathrm{X})\)）的点处的连续性。给定函数 \(f\in \mathcal{K}(\mathrm{X})\) 与数 \(\varepsilon >0\)，问题在于证明：存在 \(e\) 在 G 中的邻域 U 与 \(\mu_0\) 在 \(\mathcal{M}_+(\mathrm{X})\) 中的邻域 W，使得关系 \(s\in \mathrm{U}\)、\(\mu \in \mathrm{W}\) 蕴含

\[
\left| \int f (s x) d \mu (x) - \int f (x) d \mu_ {0} (x) \right| \leqslant \varepsilon .\tag{6}
\]

设 \(\mathbf{V}\) 为 \(\mathbf{K}\)（即 \(f\) 的支集）在 \(\mathbf{X}\) 中的一个紧邻域，并设 \(\varphi \in \mathcal{K}_{+}(\mathbf{X})\) 满足 \(\varphi(x) = 1\)（在 \(\mathbf{V}\) 上）；存在 \(\mathrm{W}_0\)（\(\mu_0\) 在 \(\mathcal{M}_{+}(\mathbf{X})\) 中的一个邻域），使得 \(a = \sup_{\mu \in \mathrm{W}_0} \mu(\mathrm{V})\) 为有限：只需对 \(\mathrm{W}_0\) 取 \(\mu \in \mathcal{M}_{+}(\mathbf{X})\) 中满足 \(|\langle \varphi, \mu - \mu_0 \rangle| \leqslant 1\) 者所成的集合即可。另一方面，由于映射 \((s, x) \mapsto sx\) 连续，存在 \(\mathrm{U}_0\)（\(e\) 在 \(\mathbf{G}\) 中的一个紧邻域），使得 \(s\mathrm{K} \subset \mathrm{V}\) 对一切 \(s \in \mathrm{U}_0\) 成立；于是函数 \((s, x) \mapsto f(sx)\) 在 \(\mathrm{U}_0 \times \mathrm{V}\) 中一致连续，故存在邻域 \(\mathrm{U} \subset \mathrm{U}_0\)（\(e\) 的一个邻域），使得 \(|f(sx) - f(x)| \leqslant \varepsilon / 2a\) 对一切 \(s \in \mathrm{U}\) 与 \(x \in \mathrm{V}\) 成立。因此，对 \(s \in \mathrm{U}\) 与 \(\mu \in \mathrm{W}_0\)，我们有

\[
\left| \int f (s x) d \mu (x) - \int f (x) d \mu (x) \right| \leqslant \varepsilon / 2;
\]

若 \(\mathbf{W} \subset \mathbf{W}_0\) 是 \(\mu_0\) 在 \(\mathcal{M}_{+}(\mathbf{X})\) 中的由测度 \(\mu \in \mathbf{W}_0\)（满足 \(|\int f(x) d\mu(x) - \int f(x) d\mu_0(x)| \leqslant \varepsilon/2\)）所组成的邻域，则 \(\mathbf{U}\) 与 \(\mathbf{W}\) 满足要求。

